\documentclass[10pt,a4paper]{amsart}
\usepackage[margin=19mm]{geometry}
\usepackage{amsmath,amssymb,mathtools}
\usepackage[scr=rsfs,cal=euler]{mathalfa}
\usepackage{xfrac}
\usepackage[unicode,hidelinks,bookmarks=false]{hyperref}
\newcommand{\ie}{i.e.\ }
\newcommand{\cf}{cf.\ }
\newcommand{\resp}{respectively\ }
\newcommand{\Subsection}{\subsection}
\providecommand{\nobreakdash}{\nobreak\hbox{-}}
\setlength{\emergencystretch}{2em}
\multlinegap0pt
\usepackage{mathtools}
\usepackage{amssymb}
\usepackage{xcolor}
\usepackage{tikz}
\usepackage{float}
\newtheorem{theorem}{Theorem}
\title{Weighted Sup-Norms at Infinity:\\ Exhaustion Functions, Decay Classes, and Anisotropic Ends}
\author{Armen Petrosyan}
\date{Source: arXiv:2512.11810v2 (CC BY 4.0)}
\begin{document}
\maketitle

\noindent\emph{Source and licence.} Weighted Sup-Norms at Infinity:\\ Exhaustion Functions, Decay Classes, and Anisotropic Ends. Version arXiv:2512.11810v2 (CC BY 4.0), distributed by arXiv under the Creative Commons Attribution 4.0 International licence, http://creativecommons.org/licenses/by/4.0/, as recorded in the arXiv licence metadata for this exact version and shown on the abstract page https://arxiv.org/abs/2512.11810v2. Full licence notice: \texttt{licenses/arxiv-2512.11810-CC-BY-4.0.txt}.\par\smallskip

\noindent\textit{Editorial scope.} Counted unit: the article's theorem thm:coarse-stability (source0.tex lines 872--882). The article's complete original proof of it is delivered with it (source0.tex lines 884--941, one \texttt{proof} environment of 58 lines). No mathematics is written, repaired or re-derived: the statement and the proof are the article's own lines, in the article's own notation, and no retained line is substituted or re-flowed.  No further source-local context is needed: every cross-reference of every retained line resolves inside the two delivered spans.  Nothing was omitted from any retained span, so the package carries no omission record.  No retained line contains a citation, so the wrapper carries no bibliography and no bibliography entry.  No retained line contains a figure environment, an image-inclusion command or drawing code, so no image is delivered and the package declares no asset dependency.  Editorial formatting only, in addition: an AMS \texttt{amsart} preamble that restates the article's own declarations and macros used by the retained lines, namely the environments \texttt{theorem} on separate counters, together with the standard packages those lines need.  The retained environments print in this document as Theorem 1 in order of appearance, whereas the article prints the same items with the numbering of its own full document.  The complete native source of the article as distributed in the arXiv e-print for this exact version is bundled unchanged as \texttt{sources/190900/source0.tex}.
\begin{theorem}[Characterisation of coarse-affine invariance]\label{thm:coarse-stability}
Let $h$ be a continuous proper exhaustion of $A$ and $\varphi\in\Phi_{\mathrm{adm}}$.
\begin{enumerate}
\item[(a)] \textup{(VAN)} holds if and only if $D_h(\lambda)<\infty$ for every $\lambda\ge1$.
\item[(b)] \textup{(INV)} holds if and only if $D_h(\lambda)<\infty$ and
$E_h(\lambda)<\infty$ for every $\lambda\ge1$.
\item[(c)] If $h(A)\supseteq[t_0,\infty)$ for some $t_0<\infty$ --- in particular if $A$ is
connected, since then $h(A)$ is an unbounded interval --- then \textup{(INV)},
\textup{(VAN)} and globally bounded dilation of $\varphi$ are all equivalent.
\end{enumerate}
\end{theorem}

\begin{proof}
Throughout, $h\sim h'$ implies that $h\to\infty$ exactly when $h'\to\infty$, so the two
$\limsup$'s are taken along the same filter. Write $W=\varphi\circ h$, $W'=\varphi\circ h'$.

\emph{(a), sufficiency.} Let $a_1h-b_1\le h'\le a_2h+b_2$. By (A2) and monotonicity,
\[
W'=\varphi(h')\le\varphi(a_2h+b_2)\le K\varphi(b_2)\,\varphi(a_2h)
\le K\varphi(b_2)\,D_h(\max(1,a_2))\,W ,
\]
using $D_h$ at the points $h(a)\in h(A)$ (for $a_2<1$ monotonicity suffices). Multiplying
by $|f-L|$ and taking $\limsup$ gives $C^{(h')}_\varphi\le C\,C^{(h)}_\varphi$, hence (VAN).

\emph{(a), necessity.} Suppose $D_h(\lambda_0)=\infty$ for some $\lambda_0\ge1$. Choose
$t_n\in h(A)$, which after passing to a subsequence we may take strictly increasing with
$t_n\to\infty$, such that $r(t_n):=\varphi(\lambda_0t_n)/\varphi(t_n)\ge n^{2}$; this is
possible because $\varphi(\lambda_0\,\cdot)/\varphi$ is bounded on compact intervals
($\varphi$ is continuous and $\ge1$), so the supremum can only be infinite along a sequence
tending to infinity. Put $h':=\lambda_0h$, a continuous proper exhaustion with $h'\sim h$.
Pick a continuous non-decreasing $v:[0,\infty)\to[1,\infty)$ with $v(t_n)=n$ and
$v(t)\to\infty$ (linear interpolation between the $t_n$, which is legitimate because they
are strictly increasing), and set
\[
f:=L+\frac{1}{\varphi(h)\,v(h)}\ \in C(A).
\]
Then $\varphi(h)|f-L|=1/v(h)\to0$, so $C^{(h)}_\varphi(f;L)=0$. Since $t_n\in h(A)$ there
is $a_n\in A$ with $h(a_n)=t_n$, and
\[
\varphi(h'(a_n))\,|f(a_n)-L|=\frac{\varphi(\lambda_0t_n)}{\varphi(t_n)\,v(t_n)}
=\frac{r(t_n)}{n}\ \ge\ n\longrightarrow\infty,
\]
so $C^{(h')}_\varphi(f;L)=\infty$ and (VAN) fails. Note that this argument uses
$t_n\in h(A)$ in an essential way: it is exactly the step that requires the bad scales to
be visible to $h$.

\emph{(b), sufficiency.} The upper bound $W'\le CW$ is as in (a). For the lower bound, let
$\lambda:=\max(1,a_1^{-1})$. On $\{h\ge b_1/a_1\}$ we have $a_1h-b_1\ge0$ and, by (A2),
$\varphi(a_1h)=\varphi\bigl((a_1h-b_1)+b_1\bigr)\le K\varphi(b_1)\varphi(a_1h-b_1)$, while
$W=\varphi(h)\le E_h(\lambda)\,\varphi(h/\lambda)\le E_h(\lambda)\varphi(a_1h)$ if
$a_1\le1$ (and $W\le\varphi(a_1h)$ if $a_1\ge1$, by monotonicity). Combining,
$W\le E_h(\lambda)K\varphi(b_1)\varphi(a_1h-b_1)\le E_h(\lambda)K\varphi(b_1)W'$. On the
complementary set $\{h<b_1/a_1\}$, which is compact, both $W$ and $W'$ are bounded above
and below by positive constants. Hence $W\asymp W'$ on $A$, and (INV) follows.

\emph{(b), necessity.} For $\mu>0$ the function $h':=\mu h$ is a continuous proper
exhaustion with $h'\sim h$. Applying (INV) to $f:=L+1/\varphi(h)$ gives
$C^{(h)}_\varphi(f;L)=1$ and
$C^{(h')}_\varphi(f;L)=\limsup\varphi(\mu h)/\varphi(h)$, so this $\limsup$ is finite and
non-zero; running the same argument with $f:=L+1/\varphi(\mu h)$ gives the reciprocal
bound. Since $\varphi(\mu\,\cdot)/\varphi$ is bounded above and below by positive constants
on compacta, we conclude $\sup_{t\in h(A)}\varphi(\mu t)/\varphi(t)<\infty$ and
$\sup_{t\in h(A)}\varphi(t)/\varphi(\mu t)<\infty$ for every $\mu>0$, i.e.\ $D_h$ and
$E_h$ are finite.

\emph{(c).} If $h(A)\supseteq[t_0,\infty)$ then $D_h(\lambda)<\infty$ forces
$D(\lambda)<\infty$, since $\sup_{t\le t_0}\varphi(\lambda t)/\varphi(t)\le\varphi(\lambda t_0)<\infty$.
Conversely $D(\lambda)<\infty$ gives both $D_h(\lambda)\le D(\lambda)$ and
$E_h(\lambda)\le D(\lambda)$. Now apply (a) and (b).
\end{proof}

\end{document}
